% Generated draft: requires mathematical and source review.
\documentclass[11pt]{article}
\usepackage[margin=25mm]{geometry}
\usepackage{fontspec}
\setmainfont{Latin Modern Roman}
\usepackage{amsmath,amssymb,mathtools,unicode-math}
\setmathfont{Latin Modern Math}
\usepackage{enumitem,needspace,xurl,hyperref}
\setlength{\parindent}{0pt}
\setlength{\parskip}{5pt}
\setlength{\emergencystretch}{3em}
\newcommand{\N}{\mathbb{N}}
\newcommand{\Z}{\mathbb{Z}}
\newcommand{\Q}{\mathbb{Q}}
\newcommand{\R}{\mathbb{R}}
\newcommand{\C}{\mathbb{C}}
\newcommand{\F}{\mathbb{F}}
\newcommand{\A}{\mathbb{A}}
\newcommand{\PP}{\mathbb{P}}
\newcommand{\E}{\mathbb{E}}
\newcommand{\eps}{\varepsilon}
\newcommand{\dd}{\,\mathrm{d}}
\begin{document}
\section*{UnsolvedMath ID: 700791}
\textbf{Source problem code:} catalogue entry
\subsection{Definitions and mathematical statement}
Let \(G\) be a finite group, let \(\pi(G)=\{p:p\text{ is prime and }p\mid |G|\}\), and let \(\operatorname{Irr}_{\C}(G)\) denote the set of irreducible complex characters of \(G\). For a prime \(p\) and a positive integer \(n\), write \(v_p(n)=\max\{a\in\Z_{\geq 0}:p^a\mid n\}\). For \(\chi\in\operatorname{Irr}_{\C}(G)\), its \(p\)-defect is
\[
 d_p(\chi)=v_p(|G|)-v_p(\chi(1)).
\]
Thus \(\chi\) has \(p\)-defect zero precisely when \(v_p(\chi(1))=v_p(|G|)\). Define
\[
 D_G(\chi)=\{p\in\pi(G):d_p(\chi)=0\},\qquad
 X(G)=\{\chi\in\operatorname{Irr}_{\C}(G):|D_G(\chi)|\geq 2\}.
\]
Under the reading that each character in question must have defect zero simultaneously for at least two distinct primes, the problem is to find necessary and sufficient conditions, expressed in group-theoretic terms, for
\[
 X(G)\neq\varnothing,
\]
and to express \(|X(G)|\) in the same terms. Each irreducible character is counted once, irrespective of how many primes belong to \(D_G(\chi)\).

The supplied wording does not explicitly settle whether the characters must be the same for the different primes. An alternative reading asks for group-theoretic necessary and sufficient conditions for
\[
 \left|\{p\in\pi(G):\exists\chi\in\operatorname{Irr}_{\C}(G)\text{ with }d_p(\chi)=0\}\right|\geq 2,
\]
allowing different characters for different primes. On that reading, the phrase ``the number of such characters'' requires further clarification. The simultaneous reading above is therefore a statement-only interpretation requiring review. The source also does not specify which group-theoretic data are permitted in the requested characterization and counting expression.
\subsection{Short English statement}
For a finite group, characterize in group-theoretic terms when an irreducible complex character has defect zero for at least two distinct prime divisors of the group order, and count all such characters. This uses the simultaneous interpretation of the supplied wording; whether different characters may serve for different primes remains unresolved.
\subsection{Sources}
\begin{itemize}
\item Supplied local UnsolvedMath registry record, permanent ID \(700791\), exact source code \texttt{catalogue entry}, title \emph{Kourovka Notebook Problem 11.99}; upstream revision \texttt{d371ee917ac9647fbb06cd6fc529deec6f3dda82}.
\item Attributed in the supplied record to S.\ P.\ Strunkov, proposed in 1990; original problem 11.99 in \emph{The Kourovka Notebook (Unsolved Problems in Group Theory)}, 11th issue (1990). The record identifies the 21st issue edition (2026), \texttt{arXiv:1401.0300v48}, as its source edition. Its collection description names E.\ I.\ Khukhro and V.\ D.\ Mazurov as editors.
\item The local record labels the problem \emph{open}. This is a supplied registry classification, not an independently verified mathematical or literature assessment. The cited source edition and its wording have not been independently checked for this draft.
\end{itemize}
\end{document}
